\documentclass[a4paper,openright,10pt]{book}

\usepackage[utf8]{inputenc}
%PAQUETEB para incluir acentos al escribir en Castellano.
\usepackage[spanish,es-tabla]{babel}
%PAQUETEB para escribir en castellano.
\usepackage{fancyhdr}
%PAQUETEB para definir encabezado y pie de páginas.
\usepackage{ragged2e}
%PAQUETEB utilizado para alinear, justificar o centrar el texto de la memoria.
\usepackage{setspace}
 %PAQUETEB para delimitar el interlineado del texto.
\usepackage{cite}
%PAQUETEB para citar y referencia a lo largo del texto.
\usepackage{enumerate} 
%PAQUETEB para formar listas organizadas y darles formato. 
\usepackage[font={color=RBlue},figurename=Fig.]{caption} 
%Paquete para poner subtítulos a las imágenes y concretamente personalizado para que sean de colo azul (elección personal).
\usepackage{graphicx} 
%Paquete para utilizar distintos colores tanto en la escritura, como en tablas, títulos y demás.
\usepackage{subfigure} 
%Paquete para incluir sub-figuras.
\usepackage{hyperref} 
%Paquete para incluir hyperlinks en el PDF.
\usepackage{eurosym} 
%Paquete para incluir el símbolo "€" en el texto.
\usepackage{pdfpages} 
%Paquete para incluir PDFs a lo largo de la memoria, y modificar su formato.
\usepackage{multirow, array} 
%Paquete para dar formato a las tablas.
\usepackage{float} 
%Paquete para forzar que una foto se introduzca exactamente en un sitio determinado y que LATEX no la reposicione.
\usepackage{longtable} 
%Paquete para la gestión y el formateo tablas de mayor dimensión.
\usepackage{xcolor,colortbl} 
%Paquete para crear colores específicos según RGB.
\usepackage{geometry}
%Paquete para definir las dimensiones de los márgenes.

\definecolor{RBlue}{RGB}{23,33,110} 
\definecolor{Rojo}{RGB}{255,0,0}
\definecolor{Cyan}{RGB}{214,234,240}
\definecolor{Naranja}{RGB}{255,222,199}
\definecolor{GrisTabla}{RGB}{245,245,245}

\setcounter{secnumdepth}{3} 
%Esta función permite que en el índice se indique hasta el sub-nivel número tres. Es decir que aparezcan hasta los apartados 1.1.1. Si se indicase {4}, se podrían observar hasta los sub-apartados 1.1.1.1. 

\setlength{\headsep}{0.5in} 
%La función "\setlength" permite cambiar (sobrescribir valores predeterminados) todas las distancias del documento. En función de que {parámetro} se indique se cambiará una distancia u otra. En este caso se está cambiando la distancia entre el encabezado y el texto para que esta sea durante todo el documento media pulgada. Se puede indicar la distancia que se quiera (cuestiones estéticas) y también se puede indicar la distancia en el sistema métrico (mm, cm).

\onehalfspace 
%Esta función determina que el interlineado a lo largo del documento sea de 1.5. Existen otras funciones para cambiar el interlineado como "\Doublespacing" para interlineado doble, "\singlespace" para interlineado sencillo o "\setspace{}" para indicar aquel que sea preferido de forma específica.

\setlength{\parindent}{0cm}
%Nuevamente se utiliza esta función para cambiar distancias del documento. Concretamente, en este caso se utiliza para cambiar la distancia de las sangrías al comenzar nuevo párrafo. Al marcar el valor en 0, se eliminan las sangrías del documento. Siempre se puede en momentos posteriores forzar la inclusión de alguna sangría o reescribir de nuevo el valor para indicar que a partir de ahí, comiencen las sangrías de nuevo.

\includeonly{Chapter/1.Intro, Chapter/2.Objetivos, Chapter/Funciones}
%Esta función es de gran importancia. Como se puede observar a la izquierda, hay una carpeta que se ha creado con diferentes capítulos. Estos capítulos son los diferentes apartados de la memoria en los que se va a ir escribiendo. Por ejemplo, el capítulo de la introducción, el capítulo de objetivos, el capítulo de la memoria técnica, o el capítulo del plan de trabajo. Bien, de cara a poder incluirlos a posteriori en esto documento principal (MAIN.tex) es necesario indicarle a LATEX que los cargue y eso se realiza con esta función. Así, habrá que incluir tantos parámetros como capítulos se deseen cargar. Los parámetros son simplemente la dirección de la carpeta en la que se encuentran, y su nombre.

\geometry{top=3cm, bottom=3cm, left=3.5cm, right=2.5cm}
%Esta función determina los márgenes que se han de utilizar a lo largo de todo el documento. Así, el margen superior será de 3 centímetros, igual que el inferior y los márgenes exterior e interior serán de 3.5 y 2.5 centímetros respectivamente.

\pagestyle{empty}
%Esta función es bastante útil ante diversos problemas que puedan surgir a lo largo del documento con las páginas en blanco. Esta función determina que el estilo (encabezados y pies de página) de la página estrictamente posterior sea totalmente blanco (vacío - empty).

%%%%%%%%%%%%%%%%%%%%%%%%%%%%%%%%%%%%%%%%%%%%%%%%%%%%%%%%%%%%%%%%
%%%% COMIENZA EL DOCUMENTO (con la función \begin{document} %%%%
%%%%%%%%%%%%%%%%%%%%%%%%%%%%%%%%%%%%%%%%%%%%%%%%%%%%%%%%%%%%%%%%

\begin{document}

\pagestyle{fancy}
%Esta función es utilizada para crear un nuevo estilo de formato, es el formato "fancy". En él vamos a determinar a continuación cómo se quieren posicionar pies de página, encabezados y demás.

\lhead{}
%Esta función sirve para que no se escriba nada (parámetro obligatorio vacío) en el encabezado en la posición izquierda. 

\chead{}
%Esta función sirve para que no se escriba nada (parámetro obligatorio vacío) en el encabezado en la posición central. 

\rhead{}
%Esta función sirve para que no se escriba nada (parámetro obligatorio vacío) en el encabezado en la posición derecha. 

\cfoot{}
%Esta función sirve para que no se escriba nada (parámetro obligatorio vacío) en el pie de página en la posición central. 

\fancyhead[OR]{}
%Esta función sirve para que no se escriba nada (parámetro obligatorio vacío) en el las páginas impares, en el encabezado en la posición derecha.

\fancyhead[OL]{}
%Esta función sirve para que no se escriba nada (parámetro obligatorio vacío) en el las páginas impares, en el encabezado en la posición izquierda.

\fancyhead[ER]{}
%Esta función sirve para que no se escriba nada (parámetro obligatorio vacío) en el las páginas pares, en el encabezado en la posición derecha.

\fancyhead[EL]{}
%Esta función sirve para que no se escriba nada (parámetro obligatorio vacío) en el las páginas pares, en el encabezado en la posición izquierda.

\fancyfoot[LE,RO]{\thepage} 
%Esta función sirve para indicar que el número de la página correspondiente se escriba a la izquierda en las páginas pares y a la derecha en las páginas impares.

\renewcommand{\headrulewidth}{0pt} 
%Esta función sirve para renovar todo tipo de comandos o distancias en el documento. Concretamente, en este caso se indica que no exista (grosor 0pt) una linea debajo del encabezado que lo separe del texto.

\includepdf[]{PortadaPFGVAux}
%Esta función incluye el PDF de la portada correspondiente al TFG. Dicho documento PDF ha de estar cargado en el menú de la izquierda como se encuentra esta portada auxiliar. Es importante incluir en el parámetro obligatorio la dirección a dicho PDF con su normbre correspondiente.

\newpage
\thispagestyle{empty}

\frontmatter
%Esta función sirve para que, en las primeras páginas de la memoria (hasta que se indique lo contrario con otra función), la paginación sea en números romanos como lo indican los criterios de formato en índice, resumen y demás apartados.

%como se puede ver, todo lo expuesto arriba (en términos de paginación) es un poco redundante y bastante lioso, aún así, funciona. Ahora, no tengo ninguna duda de que se puede optimizar, eliminar funciones o sustituir por otras para que sea más "user-friendly". Estaré encantado de escuchar como hacerlo, pero de momento desconozco como mejorar el código sin que genere errores. Dicho esto, como se comentaba al principio, es cuestión de dedicarle  tiempo y comprender más a fondo las funciones. 

\setcounter{page}{3}
%Esta función permite comenzar el contador de las páginas para la numeración en el número 3, es importante puesto que los primeros números van dedicados a portada y contraportada y no se numeran. Así, la primera página numerada será la 3. 

\chapter*{Resumen}
%Esta función determina el comienzo de un nuevo capítulo. El parámetro indica el título de lo que será el capítulo. Hay dos funciones que inician un nuevo capítulo "\Chapter*" y "\Chapter", con asterisco y sin él. El asterisco se utiliza para que no sea vea: "Capítulo 1: Resumen" y únicamente se vea "Resumen" (Gustos personales).

\thispagestyle{fancy}

Internet evoluciona constantemente y la forma de interactuar de los usuarios con ella. Desde la web 2.0, el contenido principal de todas las páginas, es el generado y proporcionado por los usuarios. Muchas empresas han crecido hasta niveles desorbitados obteniendo y comercializando toda la información posible de sus usuarios.\\

Con la web 3.0, se le puede devolver a los usuarios toda su información y todo su contenido. Como en la vida real, todos tenemos una cartera donde guardamos nuestro dinero y nuestras identidades. Así mismo, en internet funciona igual, tienes tu dinero e identidad en una cartera. Utilizando librerías web3 podemos comunicarnos con su cartera para poder identificarlo. Generamos un DID (Decentralized IDentifier - DNI descentralizado), con la posibilidad de abrir su archivador personal y guardar datos y documentos en él, como el que solemos tener en casa. Para que el uso de la web 3.0 sea satisfactoria, hay que conseguir que sea muy fácil iniciar sesión. Después de esto, se le asignará una cookie al usuario para que los siguientes accesos sean automáticos. En este proyecto, se investigará una implementación de identidad distribuida y sostenible.\\

Esto se realizará priorizando la facilidad de uso y maximizando la posibilidad de implementación en todo tipo de páginas. También investigaremos la posibilidad de implementar esta identificación a todo Internet, haciendo posible que estas identidades se puedan usar desde la carpeta de salud hasta en la próxima red social de éxito. Todo esto mientras los usuarios son dueños de sus datos, compartiendo solo lo que quieren y nada más.\\

\textbf{Descriptores}\\
blockchain, IPFS, SSI, Identidad, Distribuida

\newpage
%Función para incluir un salto a una nueva página.
\thispagestyle{empty}
%Función para que dicha página no tenga ningún formato (ni número, ni encabezado).

%Así con estas dos funciones de forma consecutiva se crea una página en blanco.

\chapter*{Abstract}
%Nuevo capítulo para incluir el resumen (si se quiere también en inglés).

\thispagestyle{fancy}
Lorem ipsum dolor sit amet, consectetur adipiscing elit. Aenean semper non orci at fringilla. Etiam fermentum in diam vestibulum pellentesque. Nam volutpat, velit ut euismod mollis, ipsum erat facilisis justo, et tempor est lectus nec massa. Pellentesque habitant morbi tristique senectus et netus et malesuada fames ac turpis egestas. Sed accumsan viverra neque eu blandit. Suspendisse in fermentum felis, id iaculis mi. Quisque maximus quam lacus, non lobortis libero tincidunt at. Pellentesque habitant morbi tristique senectus et netus et malesuada fames ac turpis egestas. Maecenas lectus nibh, sagittis ut felis non, fermentum ullamcorper leo. In hac habitasse platea dictumst. Orci varius natoque penatibus et magnis dis parturient montes, nascetur ridiculus mus. Fusce est dui, dapibus ultricies ipsum quis, ultrices maximus dolor. Proin finibus, lorem vulputate maximus convallis, quam tellus posuere magna, sed maximus nunc eros quis ipsum. Cras dignissim cursus lectus ac auctor. Donec suscipit vestibulum neque, sed lobortis est rhoncus non.\\

Praesent ut neque eros. Praesent vitae augue at diam tincidunt sagittis. In cursus lorem nec neque condimentum dictum. Morbi vel tristique orci. Cras luctus tempus elit tincidunt hendrerit. Proin bibendum arcu et sapien finibus vulputate sed eget leo. Duis at bibendum massa, sit amet ornare lorem. Donec dictum finibus fringilla. Duis accumsan lectus dolor, eu maximus ligula semper vel. Proin a sem tincidunt, mollis magna eget, consequat tortor. In sodales justo et varius scelerisque. Aliquam sit amet lacus mollis, tincidunt erat vitae, lacinia sem. Maecenas vel erat sagittis, semper ex in, commodo libero.\\

Integer malesuada quis elit eu eleifend. Suspendisse non vestibulum est, a fermentum urna. Donec ligula tortor, ultrices varius nisi eget, dictum malesuada augue. Fusce lacus orci, eleifend quis luctus dictum, ultricies id turpis. Vestibulum et auctor orci. Pellentesque habitant morbi tristique senectus et netus et malesuada fames ac turpis egestas. Integer facilisis neque quis dui dapibus dictum.\\

\textbf{Descriptors}\\
Lorem, ipsum, dolor, sit, amet

\newpage
\thispagestyle{empty}

\fancypagestyle{plain}
{
}
%Se define mediante esta función el formato plain, que no tiene nada, ni encabezados ni pies de página. Realmente esta función podría ir a comienzo del documento o en cualquier lugar, es decir, se puede mover.

\tableofcontents
%Esta función permite introducir el índice de capítulos (hasta el subíndice que se haya indicado previamente).

\cleardoublepage
%Esta función permite que aunque acabe en página impar, el siguiente apartado, comience en página impar. Con este comando te aseguras que se hace bien y que además se mantiene la numeración correspondiente en números romanos. Es decir, la página que se quedaría en blanco no está en blanco sino que mantiene la numeración. En caso de no querer numeración hay que utilizar los comando ("\newpage" y "\thispagestyle{empty}).

\listoffigures
% Es la función que introduce el índice de figuras. Funciona exactamente igual que el de capítulos.
\cleardoublepage

\listoftables
%Por ultimo, se introduce el índice de tablas en caso de que las hubiera en vuestro documento. 

\mainmatter
%Esta función se utiliza para indicar que a partir de ahora la numeración pasa a ser con números estándar y no con romanos como se había venido utilizando hasta ahora.

\pagestyle{fancy}
%A continuación se vuelve a redefinir el estilo “fancy” para que se adecúe al que ha de utilizarse a lo largo de toda la memoria. El estilo es el siguiente:

\lhead{}
%Encabezado a la izquierda vacío. En realidad debe de ir el nombre del capítulo en el que se está escribiendo pero ello se gestionará dentro de cada capítulo (pues el nombre cambia). Aquí se define el formato general.

\chead{}
%Encabezado en el centro vacío.

\fancyhead[OR]{PROYECTO FIN DE GRADO}
%En el encabezado a la derecha en las páginas pares se escribe: "PROYECTO FIN DE GRADO".

\fancyhead[OL]{} 
%Encabezado a la izquierda en las páginas pares vacío. Probablemente esta función sea redundante con la función \lhead y sería omitible.

\fancyhead[ER]{}
%Encabezado en las páginas impares a la derecha vacío.

\cfoot{} 
% Pie de página en el centro vacío.

\fancyfoot[LE,RO]{\thepage} 
%En el pie de página se va a escribir el número de la página correspondiente. Se escribirá en la derecha en las páginas pares y en la izquierda en las páginas impares.

\renewcommand{\headrulewidth}{0pt} 


\chapter{Introducción}\label{Int}
%Función que crea el título de capítulo y al cual se le da el nombre deseado a través de su parámetro obligatorio. Al no tener la función el “*” se escribirá también en el título del documento las palabras “Capítulo 1: …”. Además se indica, mediante la función “\label”, la correspondiente etiqueta que lleva asociada. La etiqueta sirve para que en caso de que luego se quiera hacer referencia al capítulo se haga llamando etiqueta tal que se escribiría “La información correspondiente a dicho tema se encuentra en el capítulo \ref{Int}.”

\thispagestyle{fancy}
%Función que determina que durante este capítulo se aplique el estilo Fancy.

\fancyhead[LE]{\thechapter.NOMBRE DEL PRIMER CAPÍTULO} 
%Función que se utiliza para indicar que en las páginas impares, aparezca en el encabezado en la parte izquierda, el número del capítulo con su correspondiente nombre.

Lorem ipsum dolor sit amet, consectetur adipiscing elit. Fusce bibendum mauris metus, quis pellentesque nisl vestibulum ut. Cras finibus, tortor id mattis imperdiet, tellus risus consectetur nisi, et luctus neque nisl nec massa. Sed interdum lacus eget nisl porttitor mattis. Donec volutpat blandit tortor ut porttitor. Integer nec pulvinar sapien. Integer eget odio feugiat, pretium tortor vel, dictum est. Proin ac eleifend augue, vitae facilisis dolor. Suspendisse at augue maximus, maximus est in, posuere dui. Nam id lorem et leo vehicula faucibus. Proin ac nunc sit amet metus ullamcorper fringilla non in quam. Sed at condimentum enim.\\
%Texto sin sentido y predeterminado que se irá utilizando a lo largo de todo el documento para simular la escritura de la memoria.

\setlength{\fboxsep}{5pt}
\begin{figure}[thbp]
\centering
\fbox{\includegraphics[width=0.6\textwidth, height=3.8cm]{Figures/LDeusto.jpeg}}
\caption{Logo de Deusto (Fuente: Universidad de Deusto) \cite{Deusto}} \label{fig:Deusto}
\end{figure}
%Todas estas funciones aparecen explicadas con detalle en el documento "Funciones.tex".

\textcolor{Rojo}{Como se puede observar en la imagen \ref{fig:Deusto}}: Cras neque purus, vulputate at neque a, rutrum elementum odio. Sed ultrices enim nulla, ac consectetur enim malesuada maximus. Nulla facilisi. Curabitur pharetra tortor nisl, vel molestie turpis commodo eget. Morbi viverra urna varius, condimentum nulla eu, semper mi. Aenean faucibus erat id felis consequat, sit amet porttitor nibh posuere. Proin imperdiet fermentum odio eget aliquet. Aenean quis auctor ante, ut consequat nunc. Maecenas ac lacinia risus. Donec finibus erat in mattis pharetra. Morbi sed metus eget magna luctus placerat quis eget velit. Praesent imperdiet velit ut magna bibendum pretium. Vestibulum fermentum nulla et suscipit tempor. Vestibulum facilisis vulputate faucibus. Vivamus et elementum tellus, at accumsan augue. Fusce elementum sem ut nunc efficitur ullamcorper. Sed ac erat quis massa pellentesque gravida id sit amet nulla. Nunc ultricies porttitor metus at fermentum \cite{Deusto}.\\

\setlength{\fboxsep}{0pt}
%"\Fbox" es una función que como se verá posteriormente es utilizada para recuadrar fotos o enmarcarlas con un borde negro. Esta función indica que no haya separación entre la foto y el marco, es decir, que justo aparezca el marco cuando acabe la foto. Se puede cambiar el "0pt" por otro número y se dejará un espacio alrededor de la foto hasta que comience el marco.


Cras efficitur purus ut ante sollicitudin, vel vulputate enim pulvinar. Suspendisse sit amet erat ut dui accumsan pharetra eget in arcu. Vestibulum fermentum a velit ac cursus. In tempus elit risus, a vestibulum est viverra a. Vestibulum ut nulla venenatis, congue urna quis, ullamcorper lacus. Duis ac hendrerit nisi, id imperdiet purus. Proin non iaculis sem. Suspendisse potenti. Nulla sed dui orci. Pellentesque vel feugiat quam, non eleifend purus. Donec porttitor velit vel sollicitudin cursus. Vivamus rhoncus vel risus non vestibulum. Aenean fermentum congue pretium. Donec dolor felis, iaculis iaculis gravida vitae, placerat ac enim.

\section{Sección}
Etiam interdum lectus nec elementum consequat. Maecenas at enim et ante aliquet porta pellentesque vitae libero. Integer tortor magna, efficitur in mauris vel, cursus consectetur magna. Donec nec nibh ultricies, ultricies velit interdum, porta nibh. Curabitur ornare, ex nec finibus interdum, leo felis semper dui, vitae auctor velit libero at arcu. In hendrerit tortor quis tempus efficitur. Aenean vel feugiat nisi, quis sagittis ex. Duis sit amet blandit ligula. Curabitur quis enim nibh. Nulla facilisi. Nunc nec dictum libero. Sed mattis euismod nulla et bibendum.

\begin{itemize}
\renewcommand{\labelitemi}{$\bullet$}
\setlength{\itemindent}{5mm}
    \item Curabitur ullamcorper varius congue.
    \item Vivamus eu quam sem. Aenean a ligula a est blandit dignissim vel non odio.
    \item Etiam sit amet velit quis enim porta semper sit amet vitae diam.
\end{itemize}

\subsection{Subsección}
In bibendum urna libero, ut maximus ex pharetra non. Aliquam sed metus eget lacus suscipit bibendum eget sed risus. Aenean dictum, urna eu lobortis auctor, quam sem porta ex, ut dignissim lectus sapien et dui. Phasellus eros massa, imperdiet vitae elit et, malesuada feugiat odio. Vivamus interdum turpis sit amet ligula rhoncus semper. Curabitur nec consequat libero, at suscipit neque. Donec commodo arcu vel eros feugiat, vitae hendrerit risus efficitur. Nulla convallis ex sed nisi ullamcorper feugiat.\\

\renewcommand{\arraystretch}{1.6}
\begin{table}[]
\begin{center}
\begin{tabular}{|m{7cm}| m{7cm} |}
\hline
\rowcolor{Cyan}
\centering \textbf{Lorem ipsum} & \hspace{2.75cm} \textbf{Dolor sit} \\\hline
\textbf{Consectetur adipiscing} & Elit\\ \hline
\rowcolor{GrisTabla}
\textbf{Consectetur adipiscing} & Elit \\ \hline
\textbf{Consectetur adipiscing} & Elit \\ \hline
\rowcolor{GrisTabla} 
\textbf{Consectetur adipiscing} & Elit \\ \hline
\rowcolor{Naranja} 
\textbf{In bibendum urna} & \textbf{Libero} \\ \hline
\end{tabular}
\caption{Tabla con texto por defecto (Fuente: Elaboración propia).}
\label{Medioambiente}
\end{center}
\end{table}

\subsubsection{Subsubsección}

Lorem ipsum dolor sit amet, consectetur adipiscing elit. Phasellus scelerisque sem quis sem commodo dictum. Maecenas venenatis hendrerit tortor, eget maximus dui ultrices ac. Vivamus aliquam ipsum non tellus lacinia varius. Nunc dapibus porta commodo. Praesent et porttitor nibh. Nunc consectetur congue dolor, ut venenatis leo. Suspendisse potenti. Nunc non nisi a metus vestibulum euismod. Vivamus dapibus lobortis sagittis. Fusce tincidunt neque velit, sed gravida ex interdum vel.\\

Ut commodo suscipit aliquet. Phasellus accumsan rhoncus lectus sit amet blandit. Duis nec quam et sapien blandit volutpat. Pellentesque nec nisl non tellus aliquet facilisis. Quisque dictum arcu quis leo blandit pulvinar non non nisi. Vivamus accumsan nec enim at scelerisque. Vestibulum ante ipsum primis in faucibus orci luctus et ultrices posuere cubilia Curae; Suspendisse hendrerit tellus ut massa sagittis, in pellentesque odio pulvinar. Sed tristique viverra mi, vitae ornare orci vestibulum eu. Nunc quis elit ante.\\
\newpage
\thispagestyle{empty}
\include{Chapter/2.Objetivos}
%Funciones para incluir los otros documentos que se han desarrollado. Dichos documentos son los distintos capítulos de la memoria. Toda la escritura de los mismos se realizará en los documentos de forma específica y no en este archivo “MAIN.tex”. Así que vamos a ellos.

\bibliographystyle{unsrt}
%Función que indica el estilo de bibliografía a utilizar.
\bibliography{References.bib}
%Función que incluye la biblografía en el documento.

\end{document}